{
\renewcommand{\arraystretch}{1.5}
\setlist[enumerate]{leftmargin=*, labelsep=0.5em}

\begin{tabularx}{\textwidth}{|>{\raggedright\arraybackslash}X|>{\raggedright\arraybackslash}X|}
\hline
\textbf{Tên Use case} & Duyệt bài viết \\
\hline
\textbf{ID} & UC-0040 \\
\hline
\textbf{Tác nhân} & Người dùng (bất kỳ, công khai) \\
\hline
\textbf{Mô tả} & Người dùng duyệt và đọc danh sách bài viết blog \\
\hline
\textbf{Điều kiện tiên quyết} & Không yêu cầu đăng nhập \\
\hline
\textbf{Điều kiện sau} & Danh sách bài viết hoặc nội dung bài viết được hiển thị \\
\hline
\textbf{Kích hoạt} & Người dùng mở trang blog \\
\hline
\textbf{Luồng chính} &
\begin{enumerate}
    \item Người dùng mở trang blog.
    \item Hệ thống hiển thị danh sách bài viết phân trang với: tiêu đề, tác giả, thẻ phân loại, ngày tạo.
    \item Người dùng có thể lọc theo tác giả (author\_id) và thẻ phân loại.
    \item Hệ thống hỗ trợ phân trang với số trang và giới hạn.
    \item Người dùng chọn một bài viết.
    \item Hệ thống hiển thị nội dung đầy đủ: tiêu đề, nội dung, author\_id, thẻ phân loại, created\_at, updated\_at.
\end{enumerate} \\
\hline
\textbf{Luồng thay thế} &
\begin{enumerate}
    \item \textbf{Xem bài viết cá nhân:}
    \begin{enumerate}
        \item[3A.] Người dùng đã xác thực lọc theo author\_id = bản thân để xem danh sách bài viết của mình.
    \end{enumerate}
\end{enumerate} \\
\hline
\textbf{Luồng ngoại lệ} &
\begin{enumerate}
    \item \textbf{Không tìm thấy bài viết:}
    \begin{enumerate}
        \item[2Err.] Hệ thống trả về danh sách rỗng với total\_count=0.
    \end{enumerate}
    \item \textbf{Bài viết không tồn tại:}
    \begin{enumerate}
        \item[6Err.] Hệ thống trả về lỗi 404 khi xem chi tiết bài viết không tồn tại.
    \end{enumerate}
\end{enumerate} \\
\hline
\end{tabularx}
\captionof{table}{Đặc tả Use case: UC-0040 Duyệt bài viết}
}
